\documentclass[a4paper, 10pt]{article}
\usepackage[T1, T2A]{fontenc}
\usepackage[utf8]{inputenc} 
\usepackage[english, russian]{babel}
\usepackage{amsfonts,amssymb,amsthm,mathtools,amsmath}
\usepackage{indentfirst} 
\usepackage{graphicx}
\usepackage{hyperref}
\usepackage{biblatex}
\usepackage{geometry}
\addbibresource{lit.bib}


\geometry{
	a4paper,
	total={170mm,257mm},
	left=20mm,
	top=20mm,
}
\theoremstyle{plain}
\newtheorem*{pro}{Доказательство}
\newtheorem*{ex}{Пример}
\newtheorem{lem}{Лемма}
\newtheorem{theorem}{Теорема}
\newtheorem{defin}{Определение}
\newtheorem*{rem}{Замечание}
\newtheorem*{sled}{Следствие}
\newcommand{\defeq}{:=}
\newcommand{\tang}[1][\hspace{0pt}]{\frac{\partial }{\partial #1}}
\newcommand\Lie{\operatorname{Lie}}
\newtheorem{remark}{Remark}

\newcommand{\wt}{\operatorname{wt}}


\title{$D_N$ аналог леммы 3}
\author{}
\date{}

\begin{document}
	\maketitle
	\begin{gather*}
		f_{D_N}(x,y)=x^{N-1}+xy^2, \qquad \Rightarrow \qquad
		A_0 =	\frac{\mathbb{C}[x,y]}{\left( (N-1) \cdot x^{N-2}+y^2, 2xy\right)}\\
	\end{gather*}
	используем базис \((1,x,x^2,\ldots,x^{N-2},y)\), поэтому в нуле\[
	xy=0,\qquad	y^2=-(N-1)x^{N-2},\qquad x^{N-1}=0.\]
Пусть даны однородные координаты на \(M = (s_0,s_1,\ldots,s_{N-2},\tau)\), c  соответствующими векторными полями
	\begin{gather*}
		e_0,e_1,\ldots,e_{N-2},h, \qquad	e_r=\frac{\partial}{\partial s_r},\qquad	h=\frac{\partial}{\partial \tau},\\
		e_r\longleftrightarrow x^r, \qquad	h\longleftrightarrow y.
	\end{gather*}
	В частности, 
	\begin{equation}
		e_a\circ e_b=e_{a+b}
		\qquad\text{для }a+b\le N-2.
		\label{eq:low_product}
	\end{equation}
Пусть есть фундаментальное соотношение редукции
	\begin{equation}
		e_1^{\circ(N-1)}
		= P(e_0, \dots, e_{N-2}, h) = 
		\sum_{k=0}^{N-2}\alpha_k(s,\tau)e_k
		+\beta(s,\tau)h, \quad \text{где } \quad \alpha_k(0)=0,\qquad\beta(0)=0.
		\label{eq:fundamental_relation} 
	\end{equation}
	\(
			\wt(x)=\frac1{N-1}, \wt(y)=\frac{N-2}{2(N-1)}
		\Rightarrow
		\wt(e_r)=\frac r{N-1}, \wt(h)=\frac{N-2}{2(N-1)},
		\wt(s_r)=1-\frac r{N-1}, \wt(\tau)=\frac{N}{2(N-1)}.
	\)
	Так как  \eqref{eq:fundamental_relation} имеет вес  $1$, получаем
	\(\wt(\alpha_k)=1-\frac{k}{N-1}=\wt(s_k).\)
	Следовательно, для $r,k\in\{1,\ldots,N-2\}$, верно
	\begin{equation}
		\left.\frac{\partial\alpha_k}{\partial s_r}\right|_0=0 \qquad\text{если }k\ne r.
		\label{eq:weight_diagonal}
	\end{equation}
	\begin{equation}
		\begin{aligned}
			&\text{Определим }B_r:=
			\left.\frac{\partial\alpha_r}{\partial s_r}\right|_0,
			\qquad r=1,\ldots,N-2.
			\label{eq:Br_def}
		\end{aligned}
	\end{equation}
	
	\begin{lem}[$D_N$]
		\[B_r=rB_1,\text{для } r=1,\ldots,N-2.
		\text{ Эквивалентно, } B_1:B_2:\cdots:B_{N-2}=1:2:\cdots:(N-2).\]
	\end{lem}
	\begin{proof}
		Пусть $m_s$ тензор умножения в точке $s=(s_0,\ldots,s_{N-2},\tau)$ и
		\[D_r:=\left.\frac{\partial m_s}{\partial s_r}\right|_{s=0}=\bigl(\Lie_{e_r}(\circ)\bigr)_0.\]
		$D_r$ симмертческое билинейное отображение \(D_r:T_0M\times T_0M\longrightarrow T_0M.\)
		Выберем $p,q\ge 1$ т.ч.\(p+q=N-1\), из ассоциативности и \eqref{eq:fundamental_relation},
		\[e_p\circ e_q=e_1^{\circ(N-1)}
		=
		\sum_{k=0}^{N-2}\alpha_k(s,\tau)e_k
		+\beta(s,\tau)h.\]
		Дифференцируя по $s_r$ в нуле и используя  \eqref{eq:weight_diagonal}, получаем
		\begin{equation}
			D_r(e_p,e_q)=B_r e_r+\lambda_r h,
			\label{eq:Dr_boundary}
		\quad \text{где }\lambda_r:=
		\left.\frac{\partial\beta}{\partial s_r}\right|_0.
	\end{equation}
	Теперь, для нечетного \(N\), все \(\lambda_r = 0\) из градуировки,  для четного же \(N\), только для одного \(r\) возможно $\lambda_r = \left.\frac{\partial\beta}{\partial s_r}\right|_0\ne0$ \par
		Пусть теперь $a,b\ge1$ и $a+b\le N-2.$ Из \eqref{eq:low_product},
		\(e_a\circ e_b=e_{a+b}.\)
		Применим условие интегрируемости
		\[\Lie_{X\circ Y}(\circ)
		=
		X\circ\Lie_Y(\circ)+Y\circ\Lie_X(\circ)\]
		с $X=e_a$ и $Y=e_b$, и ограничим в ноль. Получаем
		\begin{equation}
			\begin{aligned}
			D_{a+b}(U,V) = \bigl(\Lie_{e_{a+b}}(U\circ V)\bigr)_0
			=&
			\bigl(e_a \circ \Lie_{e_b}(U\circ V)\bigr)_0 + \bigl(\Lie_{e_a}(U\circ V) \circ e_b\bigr)_0
			=\\=&
			e_a\circ_0D_b(U,V)
			+e_b\circ_0D_a(U,V)
			\label{eq:HM_linearized} 
			\qquad \forall U,V\in T_0M.
		\end{aligned}
		\end{equation}
		
		Возьем $U=e_p$ и $V=e_q$. По \eqref{eq:Dr_boundary}, слева получаем
		\[
		D_{a+b}(e_p,e_q)
		=B_{a+b}e_{a+b}+\lambda_{a+b}h.
		\]
		Справа,
		\begin{align*}
			e_a\circ_0D_b(e_p,e_q)
			&=e_a\circ_0(B_b e_b+\lambda_bh) = B_b e_{a+b} + \lambda_b e_a \circ_0 h,\\
			e_b\circ_0D_a(e_p,e_q)
			&=e_b\circ_0(B_a e_a+\lambda_ah) = B_a e_{a+b} + \lambda_a e_b \circ_0 h.
		\end{align*}
		Заметим, что 
		\(e_c\circ_0 h=0 \text{ для любого }c\ge1,\)
		другими словами $x^cy=0$, что следует из $2xy=0$.
		Значит, произведения, содержащие \(\lambda\) в \eqref{eq:HM_linearized} исчезают. 
		\begin{equation}
		B_{a+b} e_{a+b} = B_be_{a+b}+B_ae_{a+b} = (B_a + B_b)e_{a+b}\qquad (a,b \ge 1, a+b \le N-2) \label{eq:linear_coef}
		\end{equation}
		Сравнивая коэффициенты и делая индукцию получаем
		\[
		B_r=rB_1,\qquad r=1,\ldots,N-2.
		\]
	\end{proof}
	\begin{lem}
		Если \(N>4\), то \(\forall r \lambda_r = 0\)
	\end{lem}
	\begin{proof}
		Переписывая \eqref{eq:HM_linearized} так же как в прошлой лемме получаем
		\begin{equation}
			B_{a+b}e_{a+b}+\lambda_{a+b}h = (B_a + B_b)e_{a+b} + \lambda_b e_a \circ h + \lambda_a e_b \circ h \qquad (a,b \ge 1, a+b \le N-2)\\
		\end{equation}
		Из \(e_c\circ_0 h=0 \text{ для любого }c\ge1,\) и \eqref{eq:linear_coef} получаем 
		\[\lambda_{a+b} \circ_0 e_{a+b} = 0 \quad (a,b \ge 1, a+b \le N-2)\quad \Rightarrow \quad \lambda_r = 0 \quad\forall 1<r<N-2  \]
		Теперь, \(\lambda_1 = \dfrac{\partial \beta}{\partial s_1}\) посмотрим на весовые соотношения на \(s_1\) и \(\beta\). Для \(\lambda_1 \neq 0\) необходимо
		\begin{gather*}
				\dfrac{N-2}{N-1} = \wt{s_1} = \wt{\beta} = \dfrac{N}{2(N-1)}\\
				\Downarrow\\
				\dfrac{N-2}{2} = 1\\
				N = 4
		\end{gather*}
	\end{proof}
	\begin{theorem}
		Существует изоморфизм многообразий Дубровина-Фробениуса \(M \simeq M_{D_N}\) при \(N \geq 5\)
	\end{theorem}
	\begin{proof}
		Посмотрим на \(s_1\). В силу весовых соотношений, \(s_1\) может присутствовать только в \(\alpha_1\), 
		\[\operatorname{wt}(s_{1}) = \dfrac{N-2}{N-1} = \operatorname{wt}(\alpha_{1}) > \operatorname{wt}(\alpha_{k>1})\]
		Значит, если \(B_1=0\), то умножение не зависит от \(s_1\) вовсе, а значит производная Ли по \(e_1\) равна 0. Пользуясь условием интегрируемости, получаем что 
		\[\operatorname{Lie}_{e_2}(\circ) = \operatorname{Lie}_{e_1\circ e_1}(\circ) = e_1 \circ \operatorname{Lie}_{e_1}(\circ) + \operatorname{Lie}_{e_1}(\circ) \circ e_1 = 0\]
		Индуктивно получим, что \[\forall k; \operatorname{Lie}_{e_k}(\circ) = 0.\]
		Так как \(\lambda_k = \dfrac{\partial \beta}{\partial s_1} = 0\), либо \(\dfrac{\partial P}{\partial \tau} = 0\), либо \(\dfrac{\partial \beta(s, \tau)}{\partial \tau} \neq 0\).
	\end{proof}
\end{document}